\documentclass[../DoAn.tex]{subfiles}
\begin{document}

Trên cơ sở bài toán đã được đặt ra ở Chương \ref{section:1.1} và định hướng giải
pháp tổng quát ở phần \ref{section:1.3}, chương này đi sâu hơn vào hai khía cạnh
cốt lõi của giai đoạn phân tích: \textit{khảo sát hiện trạng} và \textit{đặc tả
yêu cầu hệ thống}. Phần khảo sát làm rõ chân dung người dùng mục tiêu, mở rộng so
sánh các giải pháp WAF hiện hành ở mức độ chi tiết hơn so với khảo sát sơ bộ ở
Chương 1, từ đó định hình kỳ vọng thực tế của các bên liên quan. Phần đặc tả tiến
hành mô hình hóa hệ thống bằng các ca sử dụng (use case) tiêu biểu, đặc tả chi
tiết các chức năng xử lý nghiệp vụ chính của WAF, và liệt kê đầy đủ tập yêu cầu
phi chức năng làm cơ sở cho quá trình thiết kế và đánh giá ở các chương sau.

% ====================================================================
\section{Khảo sát hiện trạng}
\label{section:2.1}

\subsection{Khảo sát nhu cầu người dùng}
\label{subsection:2.1.1}

Một hệ thống WAF không phục vụ một nhóm người dùng duy nhất, mà giao tiếp đồng
thời với ba nhóm có vai trò và kỳ vọng khác biệt rõ rệt. Việc xác định trước
chân dung của từng nhóm là cơ sở để thiết kế các giao diện, mức phân quyền và
tập chức năng phù hợp.

\textbf{Kỹ sư bảo mật (Security Engineer)} là nhóm tương tác trực tiếp nhất với
WAF. Họ chịu trách nhiệm vận hành tập luật, đánh giá các cảnh báo và sự kiện
phát sinh, cũng như tinh chỉnh ngưỡng phát hiện để cân bằng giữa độ nhạy và tỷ
lệ báo nhầm. Họ cần một giao diện cho phép xem nhanh các sự kiện đang xảy ra,
truy xuất nhật ký chi tiết của một request đáng ngờ (bao gồm payload gốc, các
rule đã khớp, kết quả phân tích của mô hình ML nếu có), và khả năng can thiệp
tức thì như tạm chặn IP, vô hiệu hóa hoặc bật lại một rule cụ thể mà không cần
khởi động lại dịch vụ. Đối với họ, độ tin cậy của quyết định mà WAF đưa ra
phải được giải thích được: một cảnh báo \textit{block} cần phải đi kèm bằng
chứng cụ thể (rule nào, điểm anomaly là bao nhiêu, ML verdict nói gì), nếu
không, công cụ sẽ bị nghi ngờ và dần bị bỏ qua.

\textbf{Kỹ sư vận hành hệ thống và phát triển (DevOps/SRE)} không trực tiếp
phân tích các sự cố bảo mật, nhưng chịu trách nhiệm về sức khỏe vận hành tổng
thể: tỷ lệ uptime, độ trễ thêm vào do WAF, mức tiêu thụ tài nguyên, khả năng
mở rộng theo lưu lượng tăng. Nhóm này coi WAF như một thành phần trong stack
hạ tầng và quan tâm đến những đặc tính phi chức năng: khả năng triển khai bằng
một lệnh duy nhất (Docker Compose), khả năng cấu hình thông qua tệp văn bản
được kiểm soát phiên bản (declarative configuration), khả năng quan sát hệ
thống qua metric tiêu chuẩn (Prometheus, OpenTelemetry), và khả năng thực hiện
\textit{rolling update} mà không gây gián đoạn dịch vụ. Đặc biệt, họ đặc biệt
nhạy cảm với hành vi \textit{fail-open} -- nếu một thành phần phụ trợ (ví dụ
ML service) gặp sự cố, hệ thống không được đẩy traffic hợp lệ vào trạng thái
chặn nhầm.

\textbf{Lập trình viên ứng dụng (Developer)} là nhóm gián tiếp tương tác với
WAF: họ cần biết khi nào và tại sao một yêu cầu hợp lệ của ứng dụng bị WAF từ
chối, để điều chỉnh ứng dụng hoặc bổ sung ngoại lệ phù hợp. Đây là nhóm phát
sinh kỳ vọng cụ thể về khả năng \textit{exception per-rule} (cho phép loại trừ
một số đường dẫn, header, hay tham số nhất định khỏi phạm vi áp dụng của rule)
và cơ chế kiểm thử ngược (replay request) để xác minh hành vi của WAF khi
trao đổi với đội ngũ bảo mật.

Quá trình tham khảo các bài viết kinh nghiệm vận hành ModSecurity, các báo cáo
của OWASP về tối ưu Core Rule Set, cùng với các thảo luận trong cộng đồng
DevSecOps cho thấy ba kỳ vọng dùng chung của cả ba nhóm: (1) độ chính xác cao
trên cả hai chiều -- bắt được tấn công thật và không chặn nhầm traffic hợp lệ;
(2) khả năng quan sát và giải thích (observability) chi tiết với độ trễ thấp;
và (3) chi phí vận hành thấp, không yêu cầu kỹ năng quá chuyên biệt để duy trì.
Ba kỳ vọng này trở thành kim chỉ nam xuyên suốt trong quá trình thiết kế các
thành phần của hệ thống.

\subsection{Khảo sát các hệ thống WAF hiện có}
\label{subsection:2.1.2}

Chương 1 đã đề cập ngắn gọn đến ba lớp giải pháp WAF: mã nguồn mở triển khai
tại chỗ (ModSecurity, Naxsi, Shadow Daemon), thương mại trên nền tảng đám mây
(Cloudflare WAF, AWS WAF), và các giải pháp tích hợp sẵn trong reverse proxy
(NGINX App Protect). Phần này đi sâu hơn ở góc độ so sánh đặc tính kỹ thuật và
phân tích khoảng trống mà đề tài hướng đến lấp đầy.

\textbf{ModSecurity kết hợp OWASP Core Rule Set (CRS)} là tham chiếu thực tế
trong không gian WAF mã nguồn mở. Mô hình xử lý của ModSecurity dựa trên năm
pha (connection, request headers, request body, response headers, response
body), trong đó CRS đóng góp khoảng hơn 200 rule chia thành các nhóm tương ứng
với các nhóm tấn công OWASP. Điểm mạnh nổi bật là độ phổ rộng và tính được kiểm
chứng qua thời gian. Tuy nhiên, qua khảo sát các bài viết vận hành thực tế và
trải nghiệm cá nhân khi thử nghiệm trên ứng dụng DVWA, có thể nhận diện một số
hạn chế cụ thể: bộ luật mặc định tạo ra tỷ lệ false positive đáng kể trên ứng
dụng có payload JSON phức tạp, đòi hỏi giai đoạn tinh chỉnh (\textit{learning
mode}) kéo dài; kiến trúc module nhúng vào web server (Apache, NGINX) làm cho
việc reload luật mà không restart dịch vụ tương đối phức tạp; và ModSecurity
hoàn toàn không có giao diện quản trị tích hợp, buộc đội ngũ bảo mật phải xây
dựng lớp giao diện riêng hoặc tích hợp với các nền tảng SIEM bên thứ ba.

\textbf{Cloudflare WAF và AWS WAF} đại diện cho hướng tiếp cận \textit{managed
service}. Hai giải pháp này giải quyết tốt bài toán cập nhật luật và mở rộng
theo tải, đồng thời tích hợp sẵn các tính năng phụ trợ như chống DDoS volumetric
và mạng phân phối nội dung. Đặc biệt, Cloudflare cung cấp các ML-based rule học
từ lưu lượng toàn cầu, một lợi thế khó có với giải pháp on-premise. Nhược điểm
chính là chúng yêu cầu định tuyến toàn bộ lưu lượng ứng dụng qua hạ tầng nhà
cung cấp -- một ràng buộc không tương thích với các tổ chức có yêu cầu kiểm
soát dữ liệu nghiêm ngặt hoặc các môi trường ngắt mạng (air-gapped). Bên cạnh
đó, mô hình tính phí theo lưu lượng và rule khiến chi phí trở nên khó dự đoán,
và bản chất ``hộp đen'' khiến đội ngũ bảo mật mất khả năng kiểm tra logic phát
hiện ở mức chi tiết.

\textbf{Các giải pháp ít phổ biến hơn} bao gồm Naxsi (module Nginx, dựa trên
whitelist với rule trắng-ký-tự nghiêm ngặt), Shadow Daemon (kiến trúc daemon
độc lập, không được bảo trì tích cực sau 2020), và một số module thương mại của
F5 (NGINX App Protect) hay Imperva. Các giải pháp này hoặc thiếu giao diện quản
trị, hoặc bộ luật không phong phú, hoặc bị ràng buộc bởi giấy phép thương mại
và không đáp ứng tiêu chí mã nguồn mở của đề tài.

Bảng \ref{table:waf-comparison} tổng hợp so sánh các tiêu chí kỹ thuật quan
trọng nhất giữa các giải pháp khảo sát và hệ thống do đề tài đề xuất.

\begin{table}[h!]
\centering
\caption{So sánh kỹ thuật các giải pháp WAF khảo sát}
\label{table:waf-comparison}
\small
\begin{tabular}{|p{3.0cm}|p{1.8cm}|p{2.5cm}|p{2.3cm}|p{2.3cm}|}
\hline
\textbf{Tiêu chí} & \textbf{ModSecurity + CRS} & \textbf{Cloudflare / AWS WAF} & \textbf{Naxsi / Shadow Daemon} & \textbf{Hệ thống đề xuất} \\
\hline
Triển khai on-premise & Có & Không & Có & Có \\
\hline
Mã nguồn mở & Có & Không & Có & Có \\
\hline
Ngôn ngữ thực thi & C (module) & Closed & C (module) & Go \\
\hline
Tích hợp ML & Không & Có (đóng) & Không & Có (DistilBERT, gray-zone) \\
\hline
Anomaly scoring & Có & Một phần & Không & Có (đa lớp) \\
\hline
Dashboard tích hợp & Không & Có & Không & Có (web UI native) \\
\hline
RBAC nội bộ & Không & Có & Không & Có (3 vai trò) \\
\hline
Cấu hình rule động & Hạn chế & Có & Hạn chế & Có (hot-reload) \\
\hline
Cảnh báo đa kênh & Qua bên thứ ba & Có & Không & Có (Slack/SMTP/Webhook) \\
\hline
\end{tabular}
\end{table}

Từ so sánh trên có thể khẳng định rõ hơn khoảng trống mà đề tài nhắm tới: một
WAF mã nguồn mở, triển khai tại chỗ, tích hợp đầy đủ cả ba thành phần thường
thấy rời rạc trong thực tiễn -- engine luật anomaly scoring, lớp ML kiểm chứng
vùng xám, và một mặt phẳng quản trị (control plane) thống nhất qua dashboard
web cùng cơ chế xác thực và phân quyền nội bộ.

% ====================================================================
\section{Tổng quan chức năng}
\label{section:2.2}

\subsection{Biểu đồ ca sử dụng tổng quát}
\label{subsection:2.2.1}

Hệ thống WAF tương tác với bốn nhóm actor: \textbf{End User} -- người dùng cuối
gửi request hợp lệ đến ứng dụng đích; \textbf{Attacker} -- chủ thể có chủ đích
khai thác lỗ hổng; \textbf{Admin} -- quản trị viên cấp cao có toàn quyền cấu
hình hệ thống và quản trị tài khoản người dùng nội bộ; và \textbf{ML Service}
-- một actor hệ thống (system actor) đại diện cho microservice học máy bên
ngoài. Ngoài ra, hệ thống còn có \textbf{Editor} -- vai trò có quyền thay đổi
cấu hình WAF (rule, IP list, settings) nhưng không quản trị tài khoản người
khác; và \textbf{Viewer} -- vai trò chỉ đọc dashboard và nhật ký.

Hình \ref{fig:usecase-overview} mô tả biểu đồ ca sử dụng tổng quát. Các ca
sử dụng được phân thành bốn cụm: (i) cụm \textit{xử lý lưu lượng} (request
filtering, ML gray-zone confirm, rate limiting); (ii) cụm \textit{quản lý
cấu hình bảo mật} (rule management, IP management, system settings); (iii)
cụm \textit{quan sát hệ thống} (dashboard real-time, audit log query,
security alert); và (iv) cụm \textit{quản trị nội bộ} (authentication,
account settings, user management). Để tránh các đường nối chồng chéo,
biểu đồ chia actor thành hai phía: một phía là các tác nhân tham gia
đường nhanh xử lý request (End User, Attacker, và ML Service đóng vai trò
actor hệ thống), phía còn lại là ba vai trò quản trị (Viewer, Editor,
Admin) được tổ chức theo quan hệ kế thừa (Admin $\to$ Editor $\to$
Viewer) nhằm giảm số đường nối phải vẽ và làm nổi bật cấu trúc phân quyền
RBAC.

\begin{figure}[h!]
\centering
\includegraphics[width=0.95\textwidth]{Hinhve/uc_general.png}
\caption{Biểu đồ ca sử dụng tổng quát của hệ thống WAF }
\label{fig:usecase-overview}
\end{figure}

Các ca sử dụng tiêu biểu được đặc tả chi tiết trong các tiểu mục tiếp theo. Đối
với mỗi ca sử dụng, đề tài trình bày bốn yếu tố cốt lõi: \textit{actor} chính
và phụ, \textit{tiền điều kiện} (precondition), \textit{luồng chính} (main flow)
và \textit{luồng thay thế} (alternative flow), nhằm phục vụ trực tiếp giai đoạn
thiết kế ở Chương 4.

\subsection{Ca sử dụng: Kiểm tra và lọc request HTTP}
\label{subsection:2.2.2}

\textbf{Actor chính:} End User (vô tình) và Attacker (cố ý) -- cả hai không biết
sự tồn tại của WAF do hệ thống vận hành ở chế độ reverse proxy trong suốt.

\textbf{Tiền điều kiện:} WAF đã được khởi động, kết nối tới ứng dụng backend đã
hoạt động, các rule đã được nạp vào engine.

\textbf{Luồng chính:} Khi một request HTTP đến WAF, hệ thống lần lượt thực hiện
các bước sau. Đầu tiên, parser bóc tách bốn thành phần của request: phương thức,
đường dẫn, các tham số query, header và body. Tiếp theo, normalizer áp dụng
chuỗi biến đổi tiền xử lý (URL decode, lower-casing, compress whitespace) để
chuẩn hóa văn bản trước khi so khớp -- bước này quan trọng để chống các kỹ thuật
làm mờ. Sau đó, hệ thống kiểm tra IP nguồn: nếu thuộc blacklist động, request
bị từ chối ngay lập tức; nếu thuộc whitelist, request được bỏ qua phần lớn
kiểm tra. Tiếp đến, bộ rate limiter kiểm tra số lượng token còn lại trong bucket
của IP nguồn -- nếu hết token, request bị giới hạn. Cuối cùng, rule engine duyệt
tuần tự qua 78 rule, mỗi rule khớp đóng góp một giá trị vào tổng điểm anomaly.
Dựa trên điểm tổng và các ngưỡng cấu hình động (log/challenge/block), hệ thống
đưa ra quyết định: ALLOW (chuyển tiếp đến backend), CHALLENGE (yêu cầu xác thực
bổ sung), hoặc BLOCK (trả về 403). Đồng thời, một bản ghi audit log được ghi
lại bất kể quyết định nào.

\textbf{Luồng thay thế:} Nếu điểm anomaly rơi vào ``vùng xám'' \([3.0,\, 5.0)\),
WAF kích hoạt nhánh ML confirm. Engine gửi canonical text của request đến ML
service (FastAPI) qua HTTP POST với timeout 800 mili giây. Nếu ML phán quyết là
một trong các nhãn tấn công với độ tin cậy cao, một giá trị delta dương được
cộng vào điểm, đẩy quyết định về phía block/challenge. Ngược lại, nếu ML phán
quyết là normal với độ tin cậy cao, delta âm được cộng vào, đẩy quyết định về
phía allow. Trong trường hợp ML service không phản hồi trong khung thời gian
cho phép hoặc trả về lỗi, hệ thống chuyển sang chế độ \textit{fail-open} -- giữ
nguyên điểm rule và quyết định dựa trên điểm gốc, không gây gián đoạn xử lý
chính.

\begin{figure}[h!]
\centering
\includegraphics[width=0.6\textwidth]{Hinhve/uc_inspect_request.png}
\caption{Biểu đồ ca sử dụng phân rã -- Kiểm tra và lọc request HTTP}
\label{fig:uc-inspect-request}
\end{figure}

\subsection{Ca sử dụng: Quản lý luật (Rule Management)}
\label{subsection:2.2.3}

\textbf{Actor chính:} Admin (toàn quyền) và Editor (được phép sửa rule).

\textbf{Tiền điều kiện:} Người dùng đã đăng nhập thành công và có token JWT hợp
lệ với role phù hợp; định danh role được middleware xác thực kiểm tra trước khi
chuyển yêu cầu đến handler.

\textbf{Luồng chính:} Người dùng mở trang Rules từ dashboard. Hệ thống hiển thị
toàn bộ 78 rule kèm metadata (ID, category, severity, anomaly score, trạng thái
bật/tắt, hit count tích lũy). Người dùng có thể lọc theo nhóm tấn công hoặc mức
độ nghiêm trọng, bật/tắt một rule cụ thể (thay đổi có hiệu lực ngay mà không
cần restart), xem chi tiết pattern matching của rule, và tải file JSON định
nghĩa rule mới qua endpoint admin để bổ sung hoặc thay thế tập luật hiện hành.

\textbf{Luồng thay thế:} Khi tải file rule mới, hệ thống thực hiện validation
strict (kiểm tra schema, biên dịch trước các regex pattern) trước khi áp dụng.
Nếu file không hợp lệ, hệ thống từ chối và giữ nguyên tập luật cũ, đồng thời
trả về thông báo lỗi chi tiết. Mọi thay đổi tập luật đều được ghi audit log với
định danh người thực hiện.

\begin{figure}[h!]
\centering
\includegraphics[width=\textwidth]{Hinhve/uc_rule_management.png}
\caption{Biểu đồ ca sử dụng phân rã -- Quản lý luật}
\label{fig:uc-rule-mgmt}
\end{figure}

\subsection{Ca sử dụng: Giám sát và xem log}
\label{subsection:2.2.4}

\textbf{Actor chính:} Admin, Editor và Viewer.

\textbf{Tiền điều kiện:} Người dùng đã đăng nhập thành công.

\textbf{Luồng chính:} Trang Dashboard hiển thị các counter thời gian thực bao
gồm requests/giây, block rate, top attack categories trong cửa sổ thời gian
gần nhất, và phân bố quyết định (ALLOW/CHALLENGE/BLOCK). Người dùng có thể
chuyển sang trang Logs để xem chi tiết các bản ghi audit, áp dụng bộ lọc theo
khoảng thời gian, IP, loại tấn công, hay quyết định. Mỗi bản ghi log cho phép
mở rộng để xem chi tiết: payload gốc đã được làm sạch (sanitized), các rule đã
khớp, điểm anomaly cuối cùng, và nếu có, phán quyết của ML cùng độ tin cậy.

\textbf{Luồng thay thế:} Khi cơ sở dữ liệu chứa hàng triệu bản ghi, hệ thống áp
dụng phân trang phía server và đánh chỉ mục theo timestamp và IP để duy trì
thời gian phản hồi truy vấn dưới một giây.

\begin{figure}[h!]
\centering
\includegraphics[width=0.7\textwidth]{Hinhve/uc_monitoring.png}
\caption{Biểu đồ ca sử dụng phân rã -- Giám sát và xem log}
\label{fig:uc-monitoring}
\end{figure}

\subsection{Ca sử dụng: Quản lý IP (Blacklist / Whitelist)}
\label{subsection:2.2.5}

\textbf{Actor chính:} Admin và Editor.

\textbf{Luồng chính:} Trang IP Management cho phép thêm địa chỉ IP hoặc dải CIDR
vào blacklist (chặn hoàn toàn) hoặc whitelist (bỏ qua kiểm tra), với tùy chọn
TTL (time-to-live) cho phép tự động xóa sau một khoảng thời gian. Người dùng có
thể xem danh sách IP đang bị giới hạn tốc độ bởi rate limiter, danh sách IP bị
chặn động do vi phạm liên tục, và xóa thủ công các mục khỏi danh sách.

\textbf{Luồng thay thế:} Khi rate limiter phát hiện một IP vượt quá ngưỡng cho
phép liên tục trong một khoảng thời gian, IP đó được thêm tự động vào danh sách
chặn tạm thời với TTL mặc định cấu hình được. Người dùng admin có thể can thiệp
thủ công để gỡ chặn sớm hoặc kéo dài thời gian chặn.

\begin{figure}[h!]
\centering
\includegraphics[width=0.95\textwidth]{Hinhve/uc_ip_management.png}
\caption{Biểu đồ ca sử dụng phân rã -- Quản lý IP blacklist/whitelist}
\label{fig:uc-ip-mgmt}
\end{figure}

\subsection{Ca sử dụng: Cấu hình hệ thống}
\label{subsection:2.2.6}

\textbf{Actor chính:} Admin.

\textbf{Luồng chính:} Trang System Settings cho phép admin cấu hình các tham số
vận hành cốt lõi: URL backend đích, các ngưỡng quyết định (log/challenge/block),
endpoint và timeout của ML service, các tham số rate limiting (rate, burst),
và bật/tắt yêu cầu xác thực bắt buộc cho dashboard. Thay đổi có hiệu lực ngay
sau khi lưu, không cần khởi động lại proxy.

\textbf{Luồng thay thế:} Hệ thống validate các giá trị nhập vào trước khi áp
dụng (ví dụ ngưỡng block phải lớn hơn ngưỡng challenge, URL backend phải reachable).
Nếu validate thất bại, cấu hình cũ được giữ nguyên và trả về thông báo lỗi.

\begin{figure}[h!]
\centering
\includegraphics[width=0.95\textwidth]{Hinhve/uc_system_settings.png}
\caption{Biểu đồ ca sử dụng phân rã -- Cấu hình hệ thống}
\label{fig:uc-system-settings}
\end{figure}

\subsection{Ca sử dụng: Cảnh báo sự kiện bảo mật (Alert / Notification)}
\label{subsection:2.2.7}

\textbf{Actor chính:} Admin (cấu hình kênh và xem stats); Notifier worker
(actor hệ thống chạy nền) là thực thể thực hiện việc gửi.

\textbf{Trigger:} Decision layer phát hiện một sự kiện có mức độ nghiêm trọng
lớn hơn hoặc bằng ngưỡng cấu hình (mặc định HIGH).

\textbf{Luồng chính:} Khi engine đưa ra quyết định BLOCK hoặc CHALLENGE cho một
request, một sự kiện được đẩy không đồng bộ đến notifier qua kênh trong bộ nhớ.
Notifier worker chạy nền sẽ kiểm tra bộ lọc throttle để loại bỏ các sự kiện
trùng lặp (dedup theo cặp \texttt{(IP, ruleID)} trong cửa sổ 300 giây mặc định),
sau đó fan-out sự kiện đến các kênh đã được kích hoạt: Slack qua Incoming
Webhook, Email qua SMTP, và Generic Webhook (HTTP POST với payload JSON tùy
biến). Mỗi lần gửi được giới hạn timeout 5 giây và bộ đếm sent/failed được cập
nhật để phục vụ giám sát.

\textbf{Luồng thay thế:} Nếu một kênh fail (timeout, lỗi kết nối, mã trả về
phi-2xx), notifier ghi nhận thất bại và tiếp tục thử các kênh còn lại -- không
ngừng toàn bộ quá trình. Admin có thể gọi endpoint test notification để gửi
một sự kiện thử và nhận về kết quả từng kênh, phục vụ kiểm tra cấu hình ban đầu.

\begin{figure}[h!]
\centering
\includegraphics[width=0.95\textwidth]{Hinhve/uc_alert_notification.png}
\caption{Biểu đồ ca sử dụng phân rã -- Cảnh báo sự kiện bảo mật}
\label{fig:uc-alert}
\end{figure}

\subsection{Ca sử dụng: Xác thực và đăng ký tài khoản}
\label{subsection:2.2.8}

\textbf{Actor chính:} End user nội bộ (Security Engineer, DevOps, Developer).

\textbf{Luồng chính -- Đăng nhập:} Người dùng nhập tên đăng nhập và mật khẩu
trên trang Login. Hệ thống xác thực thông tin bằng cách so sánh hash bcrypt
(cost factor 10) lưu trong cơ sở dữ liệu. Nếu hợp lệ, hệ thống phát hành một
JWT (HS256) chứa định danh người dùng và vai trò, đặt token vào HttpOnly cookie
(chống đánh cắp qua XSS) và đồng thời cho phép sử dụng qua \textit{Bearer
header} với các API client.

\textbf{Luồng chính -- Đăng ký:} Đề tài cho phép tự đăng ký công khai để hỗ
trợ thử nghiệm và demo nội bộ, tuy nhiên endpoint công khai \texttt{POST
/register} \textbf{chỉ tạo được tài khoản với vai trò \texttt{viewer}}. Trường
\texttt{role} không được chấp nhận từ phía client để bịt lỗ leo thang đặc
quyền (privilege escalation). Việc nâng vai trò người dùng (lên \texttt{editor}
hoặc \texttt{admin}) chỉ được thực hiện qua endpoint quản trị tách biệt mô tả
ở tiểu mục \ref{subsection:2.2.10}.

\textbf{Vai trò và phạm vi quyền:} Hệ thống định nghĩa ba vai trò có cấp độ
quyền lực giảm dần.
\begin{itemize}[leftmargin=*]
    \item \textbf{\texttt{admin}} -- toàn quyền: cấu hình WAF, quản lý rule và
    IP list, quản trị tài khoản người dùng (CRUD), reset mật khẩu cho người
    khác, cấu hình kênh cảnh báo.
    \item \textbf{\texttt{editor}} -- quyền cấu hình WAF (rule, IP list, system
    settings, notification channels) nhưng \textbf{không} có quyền tạo, sửa,
    xóa tài khoản người dùng khác.
    \item \textbf{\texttt{viewer}} -- chỉ đọc dashboard, audit log, các trang
    cấu hình ở chế độ xem.
\end{itemize}

\textbf{Tiền điều kiện ngầm định:} Cơ sở dữ liệu áp đặt ràng buộc \texttt{CHECK
(role IN ('admin','editor','viewer'))} ở mức schema và \texttt{UNIQUE} trên cả
\texttt{username} lẫn \texttt{email} để loại trừ trùng lặp ngay tại tầng lưu trữ.

\begin{figure}[h!]
\centering
\includegraphics[width=0.95\textwidth]{Hinhve/uc_authentication.png}
\caption{Biểu đồ ca sử dụng phân rã -- Xác thực và đăng ký tài khoản}
\label{fig:uc-auth}
\end{figure}

\subsection{Ca sử dụng: Quản lý tài khoản cá nhân (Account Settings)}
\label{subsection:2.2.9}

\textbf{Actor chính:} Mọi người dùng đã đăng nhập, không phụ thuộc vai trò.

\textbf{Luồng chính -- Cập nhật thông tin cá nhân:} Người dùng mở trang
\texttt{settings.html}, nơi hiển thị username và role ở chế độ chỉ đọc (vì hai
trường này không phải là quyết định của bản thân người dùng) cùng email ở
chế độ chỉnh sửa được. Khi cập nhật email, hệ thống gọi \texttt{PUT
/waf-api/auth/me} với JWT của người dùng -- middleware \texttt{requireAuthN}
xác thực JWT trước khi chuyển đến handler.

\textbf{Luồng chính -- Đổi mật khẩu:} Người dùng nhập mật khẩu cũ và mật khẩu
mới (với độ dài tối thiểu 8 ký tự được kiểm tra cả ở client và server).
Endpoint \texttt{PUT /waf-api/auth/me/password} so sánh hash của mật khẩu cũ
trước khi cập nhật. Việc \textit{bắt buộc nhập lại mật khẩu cũ} đóng vai trò
phòng thủ chiều sâu: ngay cả khi attacker chiếm được JWT (ví dụ qua máy trạm
bị compromise), họ vẫn không thể thay đổi mật khẩu của nạn nhân nếu không
biết mật khẩu cũ.

\textbf{Luồng thay thế:} Nếu mật khẩu cũ không khớp, hệ thống trả mã trạng
thái 401 và không thay đổi gì. Mọi lần đổi mật khẩu thành công đều được ghi
audit log với định danh người thực hiện.

\begin{figure}[h!]
\centering
\includegraphics[width=0.85\textwidth]{Hinhve/uc_account_settings.png}
\caption{Biểu đồ ca sử dụng phân rã -- Quản lý tài khoản cá nhân}
\label{fig:uc-account-settings}
\end{figure}

\subsection{Ca sử dụng: Quản trị người dùng (Admin User Management)}
\label{subsection:2.2.10}

\textbf{Actor chính:} Admin (duy nhất).

\textbf{Tiền điều kiện:} JWT của người gọi chứa \texttt{role = "admin"};
middleware \texttt{requireAdmin} kiểm tra ràng buộc này trước khi chuyển đến
handler.

\textbf{Luồng chính:} Admin mở trang \texttt{users.html}. Hệ thống liệt kê
toàn bộ tài khoản trong cơ sở dữ liệu kèm vai trò, ngày tạo, và lần cập nhật
gần nhất. Admin có thể: \textbf{(i)} tạo người dùng mới với vai trò bất kỳ
trong ba vai trò qua \texttt{POST /waf-api/auth/users}; \textbf{(ii)} cập nhật
email hoặc thăng/giáng vai trò qua \texttt{PUT /waf-api/auth/users/\{id\}};
\textbf{(iii)} reset mật khẩu của người khác (không cần mật khẩu cũ) qua
\texttt{PUT /waf-api/auth/users/\{id\}/password}; và \textbf{(iv)} xóa người
dùng qua \texttt{DELETE /waf-api/auth/users/\{id\}}.

\textbf{Luồng thay thế -- Các bất biến (invariant) phải được bảo đảm:} Đây là
phần quan trọng nhất của ca sử dụng này, vì các thao tác quản trị tài khoản
có khả năng tạo ra trạng thái không thể phục hồi (irrecoverable state).

\begin{itemize}[leftmargin=*]
    \item \textbf{Last-admin protection:} Trước khi xóa hoặc demote bất kỳ
    tài khoản admin nào, hệ thống gọi \texttt{CountByRole("admin")} để đếm
    số admin còn lại. Nếu count = 1 và admin đó chính là target, thao tác
    bị từ chối với mã 400 và thông điệp \texttt{"cannot delete/demote the
    last admin"}. Bất biến này ngăn chặn tình huống hệ thống không còn admin
    nào (lockout vĩnh viễn).
    \item \textbf{Self-action protection:} Admin không được xóa hoặc demote
    chính tài khoản của mình. Bất biến này được kiểm tra cả ở tầng handler
    (so sánh \texttt{userID} từ JWT với \texttt{targetID} trên đường dẫn) và
    ở tầng giao diện (nút Delete bị disable cho hàng tương ứng với chính
    admin đang đăng nhập, theo nguyên tắc defense-in-depth).
    \item \textbf{Audit logging bắt buộc:} Mọi thao tác (create, update,
    reset password, delete) đều được ghi audit log với cặp \texttt{(actor,
    target, action, timestamp)} qua module \texttt{auditLogger.LogSystemEvent}.
\end{itemize}

\begin{figure}[h!]
\centering
\includegraphics[width=\textwidth]{Hinhve/uc_user_management.png}
\caption{Biểu đồ ca sử dụng phân rã -- Quản trị người dùng}
\label{fig:uc-user-mgmt}
\end{figure}

% ====================================================================
\section{Đặc tả chức năng}
\label{section:2.3}

\subsection{Đặc tả: Kiểm tra request và đưa ra quyết định}
\label{subsection:2.3.1}

\textbf{Input:} Một đối tượng \texttt{ParsedRequest} chứa năm trường đã được
chuẩn hóa: method, path, query parameters (key-value map), headers (multimap),
và body (giới hạn kích thước theo cấu hình, mặc định 10 MB).

\textbf{Xử lý:} Engine duyệt tuần tự qua tập rule đã nạp. Với mỗi rule, hệ
thống đánh giá các \textit{conditions} trên các trường liên quan của
\texttt{ParsedRequest}, áp dụng transform chain riêng của rule, và so khớp các
pattern (REGEX hoặc TOKEN với proximity). Nếu rule khớp, đóng góp một giá trị
\texttt{anomaly\_score \(\times\) severity\_multiplier} vào tổng điểm. Sau khi
duyệt hết, nếu tổng điểm rơi vào vùng xám \([3.0,\, 5.0)\), engine gọi ML hook
để điều chỉnh điểm. Cuối cùng, quyết định được đưa ra bằng cách so sánh điểm
tổng với ba ngưỡng cấu hình động.

\textbf{Output:} Một bộ ba \texttt{(decision, score, matched\_rules)} trong
đó \texttt{decision} thuộc tập \{\texttt{ALLOW}, \texttt{LOG}, \texttt{CHALLENGE},
\texttt{BLOCK}\}, đi kèm với một entry audit log đầy đủ.

\textbf{Ngưỡng quyết định mặc định:} \(\text{LOG} \ge 3.0\), \(\text{CHALLENGE}
\ge 4.0\), \(\text{BLOCK} \ge 5.0\). Các ngưỡng này có thể được điều chỉnh
trong runtime qua trang System Settings.

Quy trình tổng thể từ lúc tiếp nhận request đến khi ra quyết định và ghi log
được thể hiện trên biểu đồ hoạt động ở Hình \ref{fig:activity-pipeline}.

\begin{figure}[h!]
\centering
\includegraphics[height=0.85\textheight,keepaspectratio]{Hinhve/activity_request_pipeline.png}
\caption{Biểu đồ hoạt động -- Quy trình kiểm tra request và ra quyết định}
\label{fig:activity-pipeline}
\end{figure}

\subsection{Đặc tả: Xử lý anomaly scoring và ML gray-zone confirm}
\label{subsection:2.3.2}

\textbf{Mô hình scoring:} Điểm anomaly tổng của một request là tổng có trọng số
của các điểm thành phần do từng rule match đóng góp. Công thức tổng quát:
\[
S = \sum_{r \in R_{\text{matched}}} s_r \cdot m_r
\]
trong đó \(s_r\) là điểm gốc của rule \(r\) (định nghĩa trong file JSON) và
\(m_r\) là hệ số nhân theo mức độ nghiêm trọng (severity multiplier:
CRITICAL = 2.0, HIGH = 1.5, MEDIUM = 1.0, LOW = 0.5, INFO = 0.0).

\textbf{Vùng xám và ML confirm:} Khi \(S \in [3.0,\, 5.0)\), engine gọi
\texttt{runMLConfirm} với canonical text của request. ML service trả về cặp
\texttt{(label, confidence)}. Nếu \texttt{label} thuộc tập nhãn tấn công và
\texttt{confidence} \(\ge\) ngưỡng cấu hình (mặc định 0.8), \(S \leftarrow S
+ \Delta_+\) với \(\Delta_+ = 2.0\); nếu \texttt{label} là \texttt{normal} với
\texttt{confidence} cao, \(S \leftarrow S - \Delta_-\) với \(\Delta_- = 2.0\);
nếu không thỏa hai trường hợp trên, \(S\) giữ nguyên.

\textbf{Cơ chế fail-safe:} Gọi ML có timeout cứng 800 mili giây. Nếu timeout
hoặc service trả lỗi, engine không điều chỉnh điểm và tiếp tục quyết định dựa
trên điểm rule gốc -- đây là chính sách \textit{fail-open} đảm bảo ML không
phải là điểm hỏng đơn (single point of failure) của hệ thống.

\subsection{Đặc tả: Rate Limiting}
\label{subsection:2.3.3}

\textbf{Thuật toán:} Token Bucket per IP với hai tham số \texttt{rate}
(tokens/giây) và \texttt{burst} (kích thước bucket). Khi một request đến từ IP
nào đó, bucket của IP đó được rót đầy theo thời gian trôi qua từ lần cập nhật
trước, sau đó trừ một token. Nếu bucket trống, request bị từ chối với mã
trạng thái 429.

\textbf{Quản lý trạng thái:} Bucket được lưu trong bản đồ in-memory được bảo vệ
bởi \texttt{sync.Map} để cho phép truy cập concurrent. Một goroutine định kỳ
quét và loại bỏ các bucket không hoạt động lâu (idle TTL mặc định 10 phút) để
ngăn rò rỉ bộ nhớ với traffic có IP space lớn.

\textbf{Tích hợp:} Rate limit middleware được đặt sau IP-check middleware và
trước WAF rule engine middleware. IP nào vi phạm liên tục N lần trong cửa sổ
W được tự động bổ sung vào blacklist động với TTL ngắn (mặc định 1 giờ).

\subsection{Đặc tả: Cảnh báo sự kiện bảo mật}
\label{subsection:2.3.4}

\textbf{Các kênh được hỗ trợ:} Slack qua Incoming Webhook (payload JSON theo
chuẩn Slack), Email qua SMTP với template cấu hình được, và Generic Webhook
(HTTP POST với payload JSON do người dùng định nghĩa).

\textbf{Cấu hình per-destination:} Mỗi destination có các trường: enabled (flag
bật/tắt), URL hoặc thông tin kết nối, template payload, danh sách kênh áp
dụng severity filter, và các tham số throttle.

\textbf{Throttle:} Cơ chế dedup mặc định 300 giây -- một sự kiện với cùng cặp
\texttt{(channel, IP, ruleID)} không được gửi lặp lại trong cửa sổ này. Throttle
được hiện thực bằng một LRU cache giới hạn kích thước.

\textbf{Bộ lọc severity:} Sự kiện có severity thấp hơn \texttt{MinSeverity}
cấu hình (mặc định HIGH; thứ tự INFO \(<\) LOW \(<\) MEDIUM \(<\) HIGH \(<\)
CRITICAL) bị bỏ qua hoàn toàn.

\textbf{Tính chất async:} Hàm \texttt{notifier.Send()} trả về ngay lập tức sau
khi đẩy sự kiện vào kênh trong bộ nhớ (buffered channel). Worker goroutine
chạy nền tiêu thụ kênh này, áp dụng throttle, fan-out đến các destination, và
cập nhật bộ đếm sent/failed. Thiết kế này đảm bảo việc gửi cảnh báo không bao
giờ làm tăng độ trễ của đường nhanh xử lý request.

\textbf{Test destination:} Admin có thể gọi endpoint \texttt{POST
/notifications/test} để gửi một sự kiện thử đến một destination cụ thể và
nhận về kết quả chi tiết, hỗ trợ kiểm tra cấu hình ban đầu.

\subsection{Đặc tả: Audit Logging}
\label{subsection:2.3.5}

\textbf{Phạm vi ghi log:} Audit log lưu lại hai loại sự kiện. Thứ nhất là
\textit{request event} cho mọi request được flag (rule match với điểm \(\ge\)
LOG threshold), bao gồm các trường: timestamp, IP, method, path, score,
decision, danh sách rule ID đã khớp, ML label và confidence (nếu có gọi ML),
và một bản tóm tắt payload đã sanitize. Thứ hai là \textit{system event} cho
các thao tác quản trị nhạy cảm: đăng nhập thành công/thất bại, tạo/sửa/xóa
user, đổi mật khẩu, thay đổi cấu hình WAF, bật/tắt rule, thêm/xóa IP khỏi
blacklist.

\textbf{Lưu trữ:} Bảng \texttt{audit\_logs} trong PostgreSQL, có chỉ mục trên
\texttt{timestamp} và \texttt{ip}. Truy vấn lọc/phân trang qua các endpoint
\texttt{GET /logs} hỗ trợ tham số khoảng thời gian, IP, danh mục tấn công.

\textbf{Chính sách lưu giữ:} Mặc định giữ 90 ngày, có thể cấu hình. Một job
định kỳ (cron) xóa bản ghi quá hạn để kiểm soát kích thước bảng.

% ====================================================================
\section{Yêu cầu phi chức năng}
\label{section:2.4}

Trên cơ sở khảo sát ở mục \ref{section:2.1} và các đặc tả chức năng ở mục
\ref{section:2.3}, đề tài xác định bốn nhóm yêu cầu phi chức năng làm chuẩn
mực đo lường cho hệ thống.

\textbf{Hiệu năng.} Độ trễ thêm vào do WAF, đo ở phân vị 99 (P99) trên traffic
thông thường (không kích hoạt ML), phải dưới 5 mili giây. Khi nhánh ML được
kích hoạt, tổng độ trễ thêm vào không quá 800 mili giây nhờ timeout cứng ở
tầng client. Throughput mục tiêu đạt 1000 requests/giây trên một instance đơn
với phần cứng tiêu chuẩn (2 vCPU, 2 GB RAM).

\textbf{Độ chính xác.} Block rate trên tập 42 payload OWASP chuẩn phải đạt
\(\ge\) 95\%. Tỷ lệ false positive trên traffic hợp lệ (đo bằng cách replay
các request hợp pháp từ Juice Shop và DVWA) phải dưới 5\%. Đối với mô hình
ML, binary attack precision \(\ge\) 0.95 và recall \(\ge\) 0.95 trên tập kiểm
thử nội bộ.

\textbf{Khả dụng và độ tin cậy.} Hệ thống phải duy trì hành vi \textit{fail-open}:
nếu ML service không khả dụng, WAF không được chặn nhầm traffic hợp lệ. Nếu
cơ sở dữ liệu không khả dụng, các endpoint chỉ đọc (dashboard) có thể trả về
trạng thái degraded, nhưng đường nhanh xử lý request không được bị chặn (audit
log fallback sang file local). Hệ thống phải có khả năng \textit{hot reload}
tập luật và cấu hình mà không cần restart proxy.

\textbf{Bảo mật nội bộ.} Mật khẩu được lưu dưới dạng hash bcrypt cost factor
10. JWT được phát hành với secret tối thiểu 256 bit và lưu trong HttpOnly
cookie với cờ Secure khi triển khai HTTPS. Mọi endpoint quản trị (rule, IP
list, user management, system settings) đều được gate qua middleware
\texttt{requireAdmin}. Trang đăng ký công khai không cho phép gán role tùy
ý từ phía client (đã đóng lỗ leo thang đặc quyền). Các bất biến chống lockout
(last-admin protection) và chống tự gây hại (self-delete/self-demote
protection) được bảo đảm ở cả tầng handler lẫn tầng giao diện. Mọi thao tác
quản trị được ghi audit log.

\textbf{Khả năng triển khai và quan sát.} Toàn bộ hệ thống phải triển khai
được bằng một lệnh duy nhất \texttt{docker compose up}, không yêu cầu cấu
hình thủ công ngoài các tệp \texttt{.env} chuẩn. Hệ thống phải xuất metric
ở định dạng Prometheus qua endpoint \texttt{/metrics}, log có cấu trúc
(structured logging) dưới định dạng JSON ra stdout để dễ tích hợp với hạ
tầng logging tập trung.

\bigskip
\noindent\textbf{Kết chương.}
Chương này đã tiến hành khảo sát hiện trạng các giải pháp WAF mã nguồn mở và
thương mại, xác định ba nhóm người dùng mục tiêu với kỳ vọng cụ thể, mô hình
hóa hệ thống bằng mười ca sử dụng tiêu biểu trải đều bốn cụm chức năng (xử
lý lưu lượng, quản lý cấu hình, quan sát hệ thống, quản trị nội bộ), đặc tả
chi tiết năm chức năng nghiệp vụ cốt lõi, và liệt kê đầy đủ tập yêu cầu phi
chức năng. Đặc biệt, đề tài đã làm rõ kiến trúc xác thực và phân quyền với
RBAC ba vai trò (admin/editor/viewer), các bất biến bảo vệ chống lockout và
tự gây hại trong ca sử dụng quản trị tài khoản, cùng chính sách bịt lỗ leo
thang đặc quyền qua endpoint đăng ký công khai. Các đặc tả và ràng buộc xác
định ở chương này là đầu vào trực tiếp cho phần thiết kế và triển khai chi
tiết được trình bày trong Chương \ref{section:4.1}.

\end{document}
